\section{项目简介}

\subsection{调查背景}

近年来，环境噪声污染已成为社会关注度高、群众反映强烈的突出民生问题，本研究从不同角度，系统梳理浙江省噪声污染的治理背景与现实痛点，为后续问卷设计与调查范围的确定提供科学依据，也为精准反映杭甬温等核心城市居民的声环境感知奠定基础。

\subsubsection{政策背景：国家与浙江省噪声污染防治政策导向}

噪声污染防治事关民生福祉与生态文明建设，习近平总书记多次强调，要坚持以人民为中心的发展思想，“还自然以宁静、和谐、美丽”，打好污染防治攻坚战，提升生态环境治理现代化水平，这为新时代噪声污染治理工作指明了根本方向。在这一指引下，国家与浙江省先后出台一系列法律法规与政策文件，将噪声污染防治纳入生态文明建设与民生保障的重点工作范畴，也为本次杭甬温三地居民声环境感知调查提供了明确的政策导向、法治依据与现实背景。

《中华人民共和国噪声污染防治法》\scite{ref1}于2022年6月5日正式施行，全面强化\textbf{工业、建筑施工、交通运输、社会生活四大领域}噪声管控，明确政府、企业、公众各方责任，突出源头防控、分类管理、社会共治，着力解决群众关心的施工扰民、商业喧嚣、邻里噪声等现实难题，为全国噪声治理提供了坚实法治保障。

浙江省坚决贯彻国家部署，紧扣高质量发展建设共同富裕示范区目标，持续推进“浙里宁静”建设。先后印发《浙江省噪声污染防治行动计划（2023—2025年）》\scite{ref2}，提出到2025年全省声环境功能区\textbf{夜间达标率达到85\%以上}、噪声投诉明显下降的目标，系统推进制度完善、源头管控、执法监管与数字化治理。2026年3月1日起施行的《浙江省噪声污染防治办法》\scite{ref3}，进一步细化部门监管职责，创新推进宁静小区、宁静工地、宁静公园等创建，将噪声扰民纳入基层社会治理，推动形成齐抓共管格局。

\subsubsection{现实图景：浙江省环境噪声污染问题凸显}

随着浙江省城镇化进程持续推进、城市空间格局不断拓展与城市综合功能日趋完善，人口集聚、交通路网加密、工业生产活跃、商业活动频繁、社区生活多元化等多重因素叠加，交通干线噪声、工业企业生产噪声、建筑施工噪声、广场舞与商业经营等社会生活噪声源的数量、强度与影响范围持续扩大。环境噪声污染已不再是单一的环境问题，而是深度关联城市宜居建设、民生保障与基层社会治理的突出现实难题，长期噪声暴露不仅会损害居民听觉健康、引发焦虑烦躁、睡眠障碍等身心问题，还会降低居民生活满意度与城市归属感，成为制约城市人居环境品质提升、影响群众幸福感与获得感的重要短板。

\fig{图片1.png}{噪声污染影响传导逻辑图}

为系统把握全省声环境质量的变化规律，识别区域噪声污染突出短板，本研究依托浙江省数据开放平台公开监测数据，选取2020—2024年浙江省11个地级市连续五年的城市区域环境噪声、道路交通噪声、声环境功能区定点噪声监测结果，从平均等效连续A声级（LEQ）年度变化、昼夜声环境达标率、不同功能区噪声差异、各地级市空间分布特征以及全省整体变化趋势等多个维度，系统剖析近五年浙江省城市噪声污染现状、区域差异与演变规律，精准识别交通密集区、老旧居住区、工业园区周边等噪声敏感区域的污染特征，为后续开展居民噪声感知调研、梳理治理痛点、提出针对性治理对策提供客观的数据基础与现实依据。

（1）我们按照监测地点的行政区域代码，将三万余份监测数据按地市、年份分组，不对昼夜时间段进行区分，计算得到11个地市各年份的平均等效声级（LEQ），结果热力图如图\ref{fig:heatmap}(a)所示：

（2）为弥补平均等效声级在浙江省噪声污染程度评估时的局限性，本研究根据国家制定的《声环境质量标准（GB3096-2008）》，将监测数据根据时段、区域类型的不同进行达标情况的标注，并对数据按年份、地区进行分类，计算出了11个地级市2020年至2024年的声环境质量达标率，热力图如图\ref{fig:heatmap}(b)所示：

\begin{figure}[H]
\centering
\begin{subfigure}[b]{0.48\textwidth}
  \centering
  \includegraphics[width=\textwidth]{图片2.png}
  \caption{区域环境噪声等效声级热力图}
\end{subfigure}
\hfill
\begin{subfigure}[b]{0.48\textwidth}
  \centering
  \includegraphics[width=\textwidth]{图片3.png}
  \caption{声环境质量达标率变化热力图}
\end{subfigure}
\caption{浙江省各地级市噪声等效声级与达标率热力图（2020-2024）}
\label{fig:heatmap}
\end{figure}

从空间分布看，不同城市间噪声压力差异显著，宁波、湖州等城市声级长期处于高位，而浙西南等地区整体声级偏低，区域分化明显；从时间变化看，多数城市的噪声水平呈现阶段性波动特征，\textbf{2023年普遍回落，2024年又出现明显反弹}，反映出噪声治理仍存在反复性与不稳定性。

可见，浙江省各地级市的声环境达标率差异显著：\textbf{杭州市达标率长期处于低位，治理短板突出}；而浙西南等区达标率整体偏高，声环境质量相对稳定。时间上，多数城市的达标率呈现明显的阶段性波动，部分年份达标率出现大幅下滑，反映出噪声治理受外部因素影响较大，常态化管控体系仍需完善。

（3）考虑到噪声暴露在时间维度上并不均衡，本研究进一步将监测数据按上午、下午、晚上、夜间四个时段分别汇总，计算各时段的平均等效声级，并综合全省整体噪声污染情况绘制了双轴组合图，结果如图\ref{fig:temporal}所示：

\begin{figure}[H]
\centering
\begin{subfigure}[b]{0.48\textwidth}
  \centering
  \includegraphics[width=\textwidth]{env_noise_period.png}
  \caption{不同时段平均等效声级}
\end{subfigure}
\hfill
\begin{subfigure}[b]{0.48\textwidth}
  \centering
  \includegraphics[width=\textwidth]{env_noise_year.png}
  \caption{平均等效声级与达标率年际变化趋势}
\end{subfigure}
\caption{浙江省不同时段平均等效声级及环境噪声与声环境质量达标率变化趋势（2020-2024）\label{fig:temporal}}
\end{figure}

从时段分布看（图\ref{fig:temporal}(a)），上午、下午、晚上三个时段的平均等效声级较为接近，差异不明显，而夜间时段的平均等效声级显著低于其余三个时段，这与夜间交通流量下降、施工及商业活动基本停止、主要噪声源大幅减少的客观规律一致。需要指出的是，后文居民感知分析显示噪声干扰在夜间最为集中，说明居民对噪声的主观敏感程度并不单纯取决于声级绝对值：夜间背景声级偏低，且与居民休息时段高度重合，\textbf{同等声级造成的干扰反而更强}。因此夜间时段虽在平均声级上并非高值区间，仍需作为声环境管控的重点时段单独考量。

从年际变化看（图\ref{fig:temporal}(b)），全省平均等效声级2020-2022年维持在较高水平，2023年出现明显回落，2024年又快速回升至前期高位，体现出噪声污染防控的阶段性反复；声环境质量达标率则与声级变化呈反向联动关系，随声级的回落与反弹同步波动，说明声级水平仍是影响达标率的核心因素，整体治理成效的持续性有待加强。

\subsubsection{民生诉求：基于噪声投诉文本的痛点挖掘}

相较于宏观的监测数据，民众的真实感受更能反应实际情况。环境治理的成效最终落脚于民众的真实感受。为精准把握噪声污染的民生痛点，本研究利用Python开放环境，通过Scrapy模块和BeautifulSoup解析库，在征求意见以后，对浙江省“民呼我为”政务服务平台2023年至2026年的部分噪音投诉进行了合法定向爬取。随后，再对得到的数据进行Jieba分词处理，加入语气词等停用词进行文本过滤。最后进行词频分析并生成了词云图，直观呈现民生诉求的焦点。

\fig[0.7]{图片5.png}{浙江省噪声投诉词云图}

不难发现，“扰民”是本次分析中的高频词，清晰反映出噪声污染对居民休息、睡眠乃至身心健康的影响，“无法入睡”等表述反映了民众受噪声困扰的强烈不满，为研究提供了民意导向。从噪声来源看，“地铁”“广场舞”“夜间施工”“交通流量大”等词汇高频出现，说明\textbf{交通噪声、社会生活噪声与施工噪声}是民众噪声投诉的主要集中方向。

值得注意的是，\textbf{“投诉无果”“敷衍回复”“无法解决”}这类表述反复出现，折射出部分噪声问题存在治理回应不及时、问题解决不彻底的现实痛点，印证了从民众感知视角开展声环境治理研究的必要性。

\subsection{文献综述与研究缺口}

国内外围绕环境噪声感知与治理已积累较丰富成果。国外研究起步较早，核心结论集中在噪声感知的形成与噪声治理的路径两个方面。就噪声感知而言，居民对环境噪声的感知并非由声级单一决定，而受人口特征、居住环境、社会经济水平乃至个体噪声耐受能力（含神经生理机制）共同影响\scite{ref4}；年龄、性别、职业、文化程度等人口统计学差异均会带来感知偏差，个体耐受能力较弱的人群更易受到噪声侵扰。就噪声治理而言，国外研究侧重源头管控与效益量化，从低噪声路面、限行管控到建筑隔音优化均有涉及，并借助地理空间分析从成本—收益角度评估降噪投入与健康福祉收益\scite{ref5}。但此类结论多以发达国家城市化背景为对象，其噪声类型与治理体制与我国存在差异，相关经验需结合国情本土化适配。

国内研究近年增长较快，已在现状刻画与治理困境两方面形成共识。现状刻画层面，我国城市噪声主要来源于交通、生活与工业三类，区域差异显著，老旧小区与工业园区周边属重点敏感区域，兼具广场舞、动物吠叫等具有本土特色的生活噪声问题。治理困境层面，学者普遍指出政策落实不到位、部门监管边界模糊、责任主体不明、公众参与度不高、治理手段单一等问题，这与潘淑萍提出的监管碎片化、基层治理薄弱、群众取证维权困难等判断一致\scite{ref6}；同时亦关注到传统工业污染治理模式难以应对新型城市噪声的问题\scite{ref7}。

\fig{图片6.png}{噪声污染防治政策体系图}

综观已有研究，仍存在三方面不足。其一，\textbf{数据来源割裂}，客观监测与居民感知多被分开研究，缺乏“监测数据—投诉文本—问卷访谈”三位一体的整合分析；其二，\textbf{机制刻画不足}，感知研究以描述性统计与影响因素识别为主，对声环境满意度“如何形成”的传导路径缺少因果刻画；其三，\textbf{对策缺乏量化支撑}，治理建议趋于同质化，未对隔音设施、投诉响应机制等可干预变量的边际效果进行排序。

据此，本调查以杭州、宁波、温州三市常住居民为对象，整合监测数据、投诉文本与问卷访谈，尝试在数据整合、机制刻画、对策量化三个方面填补上述空白。

\begin{figure}[H]
\centering
\includegraphics[width=\textwidth]{研究局限与本研究创新框架.png}
\caption{研究局限与本研究创新框架}
\label{fig:research_framework}
\end{figure}

\subsection{调查目的}

本次调查旨在聚焦浙江省杭甬温三市居民声环境感知与噪声污染治理需求，系统梳理噪声污染的现实痛点与民众诉求，探究客观监测数据与居民主观体验之间的感知偏差，识别治理薄弱环节，为浙江省优化噪声污染防治策略、推进宁静城市建设提供科学依据与实践参考。

\subsection{调查意义}

\subsubsection{实践意义}

\fig{图片7.png}{噪声治理研究实践意义逻辑图}

本研究立足浙江省噪声污染防治现实需求，以杭州、宁波、温州三大核心城市为实证对象，整合环境监测数据、政务投诉文本与居民问卷调查信息，精准识别交通噪声、社会生活噪声、施工噪声等民生痛点，清晰呈现居民噪声困扰程度、影响时段与治理诉求，为浙江省噪声污染精准防控、靶向治理提供坚实民意支撑与数据参考。研究聚焦\textbf{“监测达标与居民感知脱节”}的现实矛盾，提出\textbf{政府治理、技术应用、公众参与}三维优化路径，有助于推动噪声治理从“数据达标”向“群众满意”转变，切实提升人居环境质量与居民幸福感。同时，本研究形成的主客观结合的噪声评估与治理分析模式，可为浙江省及国内同类城市推进“宁静城市”建设、落实《浙江省噪声污染防治行动计划》提供可复制、可推广的实践范式。

\subsubsection{理论意义}

本研究突破传统噪声污染研究偏重客观监测、忽视主观感知的局限，构建客观数据、主观感知、民生诉求三位一体的城市声环境研究框架，丰富与拓展了环境噪声污染防治领域的理论视角。通过交叉分析、探索性因子分析、有序Logistic回归与SEM结构方程模型等方法，系统揭示不同群体、不同区域噪声感知差异规律与声环境满意度影响机制，深化了城市居民环境感知与主观评价体系的理论认知。本研究将统计建模方法应用于民生导向的环境治理研究，形成标准化实证分析范式，为后续民生视角下的城市环境质量评价、污染治理研究提供理论参考与方法借鉴，助力完善高质量发展建设共同富裕示范区背景下的人居环境理论体系。

